% §3.1 AND Under Noise
%
% Cross-reference resolution status:
%
% EXPECTED UNRESOLVED (forward refs to stubs that resolve in later Plans):
%   \Cref{sec:three-faces}        -- chapter 1 §1.1 (will resolve in Plan 4)
%   \Cref{sec:or-not}             -- chapter 3 §3.2 (will resolve in Task 3 of Plan 3)
%   \Cref{sec:composition-theorem}-- chapter 3 §3.3 (will resolve in Task 4 of Plan 3)
%   \Cref{sec:miller-rabin}       -- chapter 3 §3.4 (will resolve in Task 5 of Plan 3)
%
% ALREADY RESOLVED (chapters 2 and 4, drafted in Plans 1 and 2):
%   \Cref{sec:bsc}                -- chapter 2 §2.1
%   \Cref{sec:axiom-1}            -- chapter 4 §4.2
%   \Cref{ax:element-indep}       -- Axiom 1 formal statement (chapter 4)
%   \Cref{sec:axiom-2}            -- chapter 4 §4.3
%   \Cref{ax:block-rates-indep}   -- Axiom 2 formal statement (chapter 4)
%   \Cref{def:bernoulli-bool}     -- chapter 4 §4.5
%
% Build warnings for the unresolved labels are expected at this stage.
% They resolve when Plan 3 Tasks 3-5 and Plan 4 land.

\section{AND Under Noise}
\label{sec:and-under-noise}

We have, by now, a clear picture of a single noisy bit: a
$\Bernoulli{\BoolType}$ instance with a false-positive rate $\fprate$ and
a false-negative rate $\fnrate$, as defined formally in
\Cref{def:bernoulli-bool}. The binary symmetric channel of \Cref{sec:bsc}
was the symmetric special case $\fprate = \fnrate$. The biased-coin
running example, introduced in \Cref{sec:three-faces}, is the general
asymmetric case.

The first question chapter 3 must answer is: what happens when we combine
two noisy bits using a Boolean operation? AND is the natural place to start,
because it is the most commonly composed gate in practice, the one behind
the Bloom filter's $k$-hash membership query, and the operation that
underlies Miller-Rabin's $k$-round test. Once we understand AND under noise,
OR and NOT follow without new machinery, as \Cref{sec:or-not} shows.

\subsection*{Setup: two independent noisy bits}

Let $\tilde{b}_1$ and $\tilde{b}_2$ be two $\Bernoulli{\BoolType}$ instances
with rates $(\fprate_1, \fnrate_1)$ and $(\fprate_2, \fnrate_2)$ respectively.
Each instance has its own channel matrix:
\[
  Q_i \;=\;
  \begin{pmatrix}
    1 - \fnrate_i & \fnrate_i \\
    \fprate_i & 1 - \fprate_i
  \end{pmatrix},
  \quad i = 1, 2.
\]
Row 1 of $Q_i$ governs positive inputs (true value $\top$); row 2 governs
negative inputs (true value $\bot$). Columns are the output labels $\top$
and $\bot$, following the convention from \Cref{sec:bsc}.

We apply the AND function to the two observed outputs. The \emph{true} value
of the AND is the AND of the two \emph{true} inputs. We want to know: what
is the probability that the AND of the noisy observations agrees with the AND
of the true values?

\subsection*{Axioms invoked: a flag before we compute}

Before proceeding to the case analysis, we must flag the independence
assumptions the derivation below will silently use. We state them
explicitly here because Part I's aim is to make these assumptions visible,
not merely convenient.

\textbf{Axiom 1} (\Cref{sec:axiom-1}, \Cref{ax:element-indep}) asserts
that error events at distinct inputs to a single predicate are mutually
independent. In the AND setting, $\tilde{b}_1$ and $\tilde{b}_2$ evaluate
two distinct input elements. Axiom 1 therefore licenses the factorization:
the event that $\tilde{b}_1$ makes an error is independent of the event that
$\tilde{b}_2$ makes an error. Every product of probabilities in the
case analysis below rests on this factorization.

\textbf{Axiom 2} (\Cref{sec:axiom-2}, \Cref{ax:block-rates-indep}) asserts
that the per-block error rates are conditionally independent of each other
given the objective set. For $\tilde{b}_1$, this means $\fprate_1$ and
$\fnrate_1$ are independent parameters; likewise for $\tilde{b}_2$. As a
consequence (\Cref{prop:stable-Q}), the channel matrix $Q_i$ is stable
across queries: the rate seen by one query of $\tilde{b}_1$ is the same rate
that applies to any other query of the same instance. This is what allows us
to read the rates $\fprate_i$ and $\fnrate_i$ as fixed constants throughout
the case analysis, rather than tracking a time-varying stochastic process.

Both axioms are in play. Axiom 1 handles the across-instance independence;
Axiom 2 handles the within-instance rate stability. Chapter 4 proves both
axioms in full (\Cref{sec:axiom-1} and \Cref{sec:axiom-2}) and shows, in
\Cref{sec:kronecker}, that the two together are exactly the content of the
Kronecker factorization theorem. What we do now is use these axioms
deliberately, not accidentally.

\subsection*{Truth table for AND: four cases}

The AND of two Boolean values is $\top$ if and only if both inputs are $\top$.
There are four combinations of true inputs; we analyze each.

\begin{center}
\begin{tabular}{cccc}
\hline
True $b_1$ & True $b_2$ & Correct AND output & $\Pr[\text{AND output} = \top]$ \\
\hline
$\top$ & $\top$ & $\top$ & $(1 - \fnrate_1)(1 - \fnrate_2)$ \\
$\top$ & $\bot$ & $\bot$ & $(1 - \fnrate_1)\fprate_2$ \\
$\bot$ & $\top$ & $\bot$ & $\fprate_1(1 - \fnrate_2)$ \\
$\bot$ & $\bot$ & $\bot$ & $\fprate_1 \fprate_2$ \\
\hline
\end{tabular}
\end{center}

The rightmost column gives $\Pr[\text{observed AND} = \top]$ for each
true-input combination. It does not give the probability that the AND is
\emph{correct}; that requires also knowing which output is correct for that
row. Let us work through each case in turn.

\medskip
\noindent\textbf{Case 1: both true inputs are $\top$.}
The correct output is $\top$. For the AND of observed outputs to equal
$\top$, both $\tilde{b}_1$ and $\tilde{b}_2$ must report $\top$. Each
does so correctly with probability $1 - \fnrate_i$ (the true-positive rate).
By Axiom 1, the two events are independent:
\[
  \Pr[\tilde{b}_1 = \top,\; \tilde{b}_2 = \top \mid b_1 = \top, b_2 = \top]
  \;=\; (1 - \fnrate_1)(1 - \fnrate_2).
\]
This is the AND's true-positive probability.

\medskip
\noindent\textbf{Case 2: $b_1 = \top$, $b_2 = \bot$.}
The correct output is $\bot$ (since $b_2 = \bot$). For the AND of observed
outputs to incorrectly equal $\top$, $\tilde{b}_1$ must report $\top$
(correctly, with probability $1 - \fnrate_1$) \emph{and} $\tilde{b}_2$ must
report $\top$ (incorrectly, a false positive, with probability $\fprate_2$).
By Axiom 1:
\[
  \Pr[\tilde{b}_1 = \top,\; \tilde{b}_2 = \top \mid b_1 = \top, b_2 = \bot]
  \;=\; (1 - \fnrate_1)\fprate_2.
\]
This is a false positive conditional on the input pattern $(\top, \bot)$.

\medskip
\noindent\textbf{Case 3: $b_1 = \bot$, $b_2 = \top$.}
Symmetric to Case 2. The AND output is $\top$ incorrectly when
$\tilde{b}_1$ makes a false positive and $\tilde{b}_2$ reports correctly:
\[
  \Pr[\tilde{b}_1 = \top,\; \tilde{b}_2 = \top \mid b_1 = \bot, b_2 = \top]
  \;=\; \fprate_1(1 - \fnrate_2).
\]

\medskip
\noindent\textbf{Case 4: both true inputs are $\bot$.}
The correct output is $\bot$. The AND output is $\top$ incorrectly only when
both $\tilde{b}_1$ and $\tilde{b}_2$ make false positives simultaneously:
\[
  \Pr[\tilde{b}_1 = \top,\; \tilde{b}_2 = \top \mid b_1 = \bot, b_2 = \bot]
  \;=\; \fprate_1\fprate_2.
\]

\subsection*{Closed-form output rates}

The AND gate's output is a new Boolean value. To characterize it as a
$\Bernoulli{\BoolType}$ instance, we need to identify its false-positive
rate and false-negative rate with respect to the \emph{true AND output}.

The true AND output is $\top$ if and only if $b_1 = b_2 = \top$. The true
AND output is $\bot$ in all three remaining cases.

\medskip
\noindent\textbf{False-negative rate of AND.}
A false negative occurs when the true output is $\top$ (meaning $b_1 = b_2 =
\top$) but the observed AND is $\bot$. From Case 1, the probability that the
observed AND is $\top$ when the true output is $\top$ is
$(1 - \fnrate_1)(1 - \fnrate_2)$. Therefore, the false-negative rate is:
\begin{equation}
  \fnrate_\wedge
  \;=\; 1 - (1-\fnrate_1)(1-\fnrate_2)
  \;=\; \fnrate_1 + \fnrate_2 - \fnrate_1\fnrate_2.
  \label{eq:and-fnr}
\end{equation}
This is the probability that at least one of the two instances reports $\bot$
when the true value is $\top$. The union-of-events formula applies because
both must succeed for the AND to succeed.

\medskip
\noindent\textbf{False-positive rate of AND.}
A false positive occurs when the true output is $\bot$ (meaning at least one
of $b_1$, $b_2$ is $\bot$) but the observed AND is $\top$. The true output
is $\bot$ for three input patterns: $(\top, \bot)$, $(\bot, \top)$, and
$(\bot, \bot)$.

The observed AND is $\top$ in each of these three cases with the probabilities
from Cases 2, 3, and 4 above. The false-positive rate of the AND gate is
therefore a mixture over these three patterns, weighted by their prior
probabilities. Let $\pi$ denote the prior probability that $b_1 = \top$, and
let $\pi_1 = \pi$, $\pi_2 = \pi$ (assuming the two bits are drawn from the
same prior for simplicity; the general formula follows with $\pi_1$ and
$\pi_2$ distinct).

The conditional-on-pattern false-positive events are:
\begin{align*}
  \Pr[\wedge = \top \mid b_1 = \top, b_2 = \bot] &= (1 - \fnrate_1)\fprate_2, \\
  \Pr[\wedge = \top \mid b_1 = \bot, b_2 = \top] &= \fprate_1(1 - \fnrate_2), \\
  \Pr[\wedge = \top \mid b_1 = \bot, b_2 = \bot] &= \fprate_1\fprate_2.
\end{align*}

These are three distinct conditional false-positive probabilities, one per
input pattern. In the Bernoulli framework, the false-positive rate of the AND
gate is the probability that the observed AND is $\top$ given that the true
AND is $\bot$. Because the three patterns that yield true AND $= \bot$ have
different likelihoods of a false $\top$ output, the effective false-positive
rate of the AND gate depends on the input distribution.

For the special case where both bits are independently queried about the same
underlying question (both $\pi_1 = \pi_2 = \pi$, and the two bits are drawn
independently), the marginal false-positive rate is:
\begin{equation}
  \fprate_\wedge(\pi)
  \;=\;
  \frac{
    \pi(1-\pi)(1-\fnrate_1)\fprate_2
    + (1-\pi)\pi\,\fprate_1(1-\fnrate_2)
    + (1-\pi)^2 \fprate_1\fprate_2
  }{
    \pi(1-\pi) + (1-\pi)\pi + (1-\pi)^2
  }.
  \label{eq:and-fpr-general}
\end{equation}
The denominator is the probability that the true AND is $\bot$, which equals
$1 - \pi^2$.

\medskip
\noindent\textbf{The input-distribution-free bounds.}
In most practical analyses, we do not know the exact prior $\pi$. What we
can say regardless of $\pi$ is the following. Since we are taking the AND
of two outputs, the AND gate can only output $\top$ if both inputs report
$\top$. This means:
\[
  \fprate_\wedge \;\leq\; \min(\fprate_1, \fprate_2),
\]
since the AND gate can falsely report $\top$ only when both individual
instances report $\top$, which is at least as unlikely as each one doing so
alone. The upper bound follows from the cases directly: no single case has a
false-positive probability exceeding $\max(\fprate_1, \fprate_2)$, and the
worst case (Case 2 with $\fnrate_1 = 0$) approaches $\fprate_2$.

The most important special case is when both instances are independent,
identically noisy copies: $\fprate_1 = \fprate_2 = \fprate$, $\fnrate_1 =
\fnrate_2 = \fnrate$. The worst-case false-positive input pattern is the
$(\bot, \bot)$ case, giving $\fprate^2$. For the $(\top, \bot)$ and
$(\bot, \top)$ patterns, the false-positive probability is
$(1 - \fnrate)\fprate$, which for small $\fnrate$ is approximately $\fprate$.
The effective false-positive rate is therefore bounded above by
$(1 - \fnrate)\fprate$, and the dominant pattern depends on the input prior.
When both bits query the same two-valued population symmetrically, the
AND gate reduces the false-positive rate dramatically compared to a single
query: the $(\bot, \bot)$ contribution goes as $\fprate^2$, which for small
$\fprate$ is negligible. This is the mechanism behind Miller-Rabin's
exponential error reduction, which we examine in \Cref{sec:miller-rabin}.

\medskip
\noindent\textbf{AND as a $\Bernoulli{\BoolType}$ operation.}
The AND gate's output is a new Boolean value. Its false-negative rate
$\fnrate_\wedge$ in~\eqref{eq:and-fnr} depends only on the input rates,
not on the input prior. Its false-positive rate $\fprate_\wedge$ depends on
the input prior through~\eqref{eq:and-fpr-general}, unless we are in a
special case that eliminates the dependence (for example, the
$(1-\fnrate)\fprate_2$ bound applies when $b_1 = \top$ is guaranteed).

The conclusion is that AND is a binary operation on $\Bernoulli{\BoolType}$
instances that produces a new $\Bernoulli{\BoolType}$ instance. The type is
closed under AND: each output is characterized by a pair of rates, derived
from the input rates via the formulas above. The derivation invoked both
axioms explicitly: Axiom 1 to factor probabilities across the two instances,
and Axiom 2 to treat each instance's rates as stable constants throughout the
computation.

\subsection*{Worked example: AND of two noisy coin observers}

The biased-coin thread from \Cref{sec:three-faces} returns here with its
most concrete appearance so far. We have two observers of the same coin flip.

\medskip
\noindent\textbf{Setup.} A coin lands heads ($\top$) with probability
$\pi = 0.5$ (fair coin). Two independent observers $\tilde{b}_1$ and
$\tilde{b}_2$ each report the outcome, but with different noise levels:
\begin{itemize}
  \item Observer 1: $\fprate_1 = 0.10$, $\fnrate_1 = 0.05$ (more likely to miss
    heads than to hallucinate heads).
  \item Observer 2: $\fprate_2 = 0.08$, $\fnrate_2 = 0.12$ (more likely to
    miss heads than to hallucinate, but with different magnitudes).
\end{itemize}
We declare ``heads'' (output $\top$) only if both observers agree: output
$\tilde{b}_1 \wedge \tilde{b}_2$.

\medskip
\noindent\textbf{False-negative rate of AND.}
From~\eqref{eq:and-fnr}:
\[
  \fnrate_\wedge
  \;=\; \fnrate_1 + \fnrate_2 - \fnrate_1\fnrate_2
  \;=\; 0.05 + 0.12 - (0.05)(0.12)
  \;=\; 0.17 - 0.006
  \;=\; 0.164.
\]
Requiring both observers to report ``heads'' when the coin actually landed
heads raises the false-negative rate from the individual rates of $5\%$ and
$12\%$ to $16.4\%$. This is the cost of the AND gate: it is harder to get
two correct $\top$ reports simultaneously than just one.

\medskip
\noindent\textbf{Conditional false-positive rates, by pattern.}
\begin{align*}
  \Pr[\wedge = \top \mid b_1 = \top, b_2 = \bot]
  &= (1 - 0.05)(0.08) = 0.95 \times 0.08 = 0.076, \\
  \Pr[\wedge = \top \mid b_1 = \bot, b_2 = \top]
  &= (0.10)(1 - 0.12) = 0.10 \times 0.88 = 0.088, \\
  \Pr[\wedge = \top \mid b_1 = \bot, b_2 = \bot]
  &= (0.10)(0.08) = 0.008.
\end{align*}
The AND gate is hardest to fool when both true values are $\bot$: the
probability of a false $\top$ output drops to $0.8\%$, compared to $7.6\%$
or $8.8\%$ when one true value is $\top$.

\medskip
\noindent\textbf{Effective false-positive rate (fair coin, $\pi = 0.5$).}
With a fair coin, the four true-input patterns have probabilities:
\begin{align*}
  \Pr[b_1 = \top, b_2 = \top] &= 0.5 \times 0.5 = 0.25, \\
  \Pr[b_1 = \top, b_2 = \bot] &= 0.5 \times 0.5 = 0.25, \\
  \Pr[b_1 = \bot, b_2 = \top] &= 0.5 \times 0.5 = 0.25, \\
  \Pr[b_1 = \bot, b_2 = \bot] &= 0.5 \times 0.5 = 0.25.
\end{align*}
The true AND is $\bot$ when any of the last three patterns hold, giving
$\Pr[\text{true AND} = \bot] = 0.75$. The marginal probability that the
observed AND is $\top$ when the true AND is $\bot$ is the weighted average:
\begin{align*}
  \text{numerator}
  &= 0.25 \times 0.076 + 0.25 \times 0.088 + 0.25 \times 0.008 \\
  &= 0.019 + 0.022 + 0.002 = 0.043.
\end{align*}
The effective false-positive rate is:
\[
  \fprate_\wedge
  \;=\; \frac{0.043}{0.75}
  \;\approx\; 0.057.
\]

\medskip
\noindent\textbf{Comparison to individual rates.}
The two individual observers have $\fprate_1 = 0.10$ and $\fprate_2 = 0.08$.
Their AND gate reduces the effective false-positive rate to approximately
$0.057$ under a fair-coin prior. The reduction is meaningful but not
dramatic, because the dominant false-positive patterns are those where one
true value is $\top$ (so only one observer needs to make an error); only the
all-negative pattern benefits from requiring both observers to err
simultaneously.

The false-negative rate rises from $5\%$ and $12\%$ individually to $16.4\%$.
The AND gate trades false-negative accuracy for false-positive suppression.
In applications where false positives are costly and false negatives are
tolerable, this is a favorable trade. In applications where false negatives
are costly (a safety system that must not miss a $\top$ event), requiring both
observers to agree is counterproductive.

\medskip
\noindent\textbf{Summary of the example.}
Two $\Bernoulli{\BoolType}$ observers with rates $(0.10, 0.05)$ and
$(0.08, 0.12)$ combine under AND to produce a new $\Bernoulli{\BoolType}$
instance with $\fnrate_\wedge = 0.164$ and $\fprate_\wedge \approx 0.057$
(at $\pi = 0.5$). The rates follow closed-form formulas derived from the
four truth-table cases. Both Axiom 1 and Axiom 2 were invoked:
Axiom 1 to factor the joint probability of both observers reporting $\top$
into a product of individual probabilities; Axiom 2 to treat each
observer's rates as stable constants rather than time-varying quantities.

\subsection*{Looking ahead}

The AND derivation gives us one gate. \Cref{sec:or-not} will derive OR and
NOT under the same methodology: NOT is elementary (it swaps $\fprate$ and
$\fnrate$), and OR follows via De Morgan's law applied to noisy gates.
Together, AND, OR, and NOT generate all of Boolean algebra.

Once we have all three gates, \Cref{sec:composition-theorem} asks what
happens when we chain them: if AND has false-positive rate $\fprate_\wedge$,
and we then AND the result with a third noisy bit, what is the total
false-positive rate? The answer is the composition theorem,
$\eta_{\text{total}} = 1 - \prod_{i=1}^k (1 - \eta_i)$, which describes
how error rates accumulate in chains of identity-with-error operations.

That composition theorem is precisely what makes Miller-Rabin's $k$-round
bound derivable from first principles, as \Cref{sec:miller-rabin} will show.
The $k$-round Miller-Rabin test is, in the framework's language, a $k$-fold
AND of $\Bernoulli{\BoolType}$ witnesses. The derivation we just completed
for two witnesses extends directly to $k$; the formulas are the same;
only the indices change.
